%%%%%%%%%%%%%%%%%%%%%%%%%%%%%%%%%%%%%%%%%%%%%%%%%%%%%%%%%%%%%%%%%%%%%
%% This is a "template" model document for submission to the
%% American Chemical Society (ACS).
%%
%% The guidance here contains information about how you may wish to
%% modify it to match the requirements of various ACS journals. The
%% ACS do *not* typeset accepted articles using LaTeX, so there is
%% no specific class required.
%%
%% This template deliberately does *not* seek to reproduce
%% the layout of the typeset journal: this is explicitly not
%% required by the ACS for LaTeX submissions.
%%
%% Please report any issues with the template at
%% https://github.com/josephwright/acs-template/issues
%%
%% Released under the Creative Commons 0 license
%% https://creativecommons.org/public-domain/cc0/
%% 
%% Copyight (c) 2025 Joseph Wright
%%%%%%%%%%%%%%%%%%%%%%%%%%%%%%%%%%%%%%%%%%%%%%%%%%%%%%%%%%%%%%%%%%%%%
\documentclass[letterpaper]{article}

%%%%%%%%%%%%%%%%%%%%%%%%%%%%%%%%%%%%%%%%%%%%%%%%%%%%%%%%%%%%%%%%%%%%%
%% Font setup - delete if you are using LuaLaTeX
%%%%%%%%%%%%%%%%%%%%%%%%%%%%%%%%%%%%%%%%%%%%%%%%%%%%%%%%%%%%%%%%%%%%%
\usepackage[T1]{fontenc}

%%%%%%%%%%%%%%%%%%%%%%%%%%%%%%%%%%%%%%%%%%%%%%%%%%%%%%%%%%%%%%%%%%%%%
%% Adjust the margins and allow for line spacing
%%%%%%%%%%%%%%%%%%%%%%%%%%%%%%%%%%%%%%%%%%%%%%%%%%%%%%%%%%%%%%%%%%%%%
\usepackage{geometry}
\geometry{margin = 1in}
\usepackage{setspace}

%%%%%%%%%%%%%%%%%%%%%%%%%%%%%%%%%%%%%%%%%%%%%%%%%%%%%%%%%%%%%%%%%%%%%
%% Reference support
%%
%% The recommended method for producing the reference section is
%% to use biblatex. If you wish to use a classical BibTeX
%% approach, this is easiest to achieve using the achemso package.
%% In that case, you should remove the biblatex lines.
%%%%%%%%%%%%%%%%%%%%%%%%%%%%%%%%%%%%%%%%%%%%%%%%%%%%%%%%%%%%%%%%%%%%%
% You can adjust the printing of DOI, article title, etc. using
% package options, e.g. "doi = true"; see the biblatex manual for
% details of adjusting the number of authors printed, e.g.
% "maxnames = 15" to print no more than 15 authors.
\usepackage[style = chem-acs]{biblatex}
\addbibresource{quantum.bib}
% If you are using classical BibTeX, remove the above lines and 
% uncomment:
%\usepackage{achemso}

%%%%%%%%%%%%%%%%%%%%%%%%%%%%%%%%%%%%%%%%%%%%%%%%%%%%%%%%%%%%%%%%%%%%%
%% Graphic inclusion and scheme and chart support
%%%%%%%%%%%%%%%%%%%%%%%%%%%%%%%%%%%%%%%%%%%%%%%%%%%%%%%%%%%%%%%%%%%%%
\usepackage{graphicx}
\usepackage{float}
\newfloat{scheme}{htbp}{los}
\floatname{scheme}{Scheme}
\floatname{chart}{Chart}
\newfloat{graph}{htbp}{loh}

%%%%%%%%%%%%%%%%%%%%%%%%%%%%%%%%%%%%%%%%%%%%%%%%%%%%%%%%%%%%%%%%%%%%%
%% Common support packages
%%%%%%%%%%%%%%%%%%%%%%%%%%%%%%%%%%%%%%%%%%%%%%%%%%%%%%%%%%%%%%%%%%%%%
%\usepackage{chemformula} % Formulas using \ch{}
% or
\usepackage[version = 4]{mhchem} % Formulas using \ce{}

%%%%%%%%%%%%%%%%%%%%%%%%%%%%%%%%%%%%%%%%%%%%%%%%%%%%%%%%%%%%%%%%%%%%%
%% Many journals require that sections are unnumbered: this 
%% is activated here
%%%%%%%%%%%%%%%%%%%%%%%%%%%%%%%%%%%%%%%%%%%%%%%%%%%%%%%%%%%%%%%%%%%%%
% \setcounter{secnumdepth}{-1}

%%%%%%%%%%%%%%%%%%%%%%%%%%%%%%%%%%%%%%%%%%%%%%%%%%%%%%%%%%%%%%%%%%%%%
%% Place any additional macros here.  Please use \newcommand* where
%% possible, and avoid layout-changing macros (which are not used
%% when typesetting).
%%%%%%%%%%%%%%%%%%%%%%%%%%%%%%%%%%%%%%%%%%%%%%%%%%%%%%%%%%%%%%%%%%%%%
\usepackage{subfiles}
\usepackage{amsmath, amssymb, braket} % AMS math
\usepackage{lmodern} % Font for math
\usepackage{bm} % Bold math
\usepackage{subcaption} % Captions for figures
\usepackage{nameref} % For cross-referencing sections
\usepackage{xr}
\externaldocument{jctc-draft_si}

\newcommand*\mycommand[1]{\texttt{\emph{#1}}}
\usepackage{jctc-draft-commands}

%%%%%%%%%%%%%%%%%%%%%%%%%%%%%%%%%%%%%%%%%%%%%%%%%%%%%%%%%%%%%%%%%%%%%
%% Author and title data:
%% the authblk package is currently the simplest way to provide this
%%%%%%%%%%%%%%%%%%%%%%%%%%%%%%%%%%%%%%%%%%%%%%%%%%%%%%%%%%%%%%%%%%%%%
\usepackage{authblk}
\author[1,3]{Hideo Takahashi*}
\affil[1]{Managed \& Platform Services Business Division, Hitachi Ltd.,
         292, Yoshida, Totsuka-ku, Yokohama, 244-0817, Japan}
\author[2,3]{Tatsuya Tomaru}
\affil[2]{Center for Exploratory Research, Research \& Development Group, Hitachi Ltd.,
Kokubunji, Tokyo, 185-8601, Japan}
\affil[3]{Institute of Industrial Science, The University of Tokyo, 4-6-1 Komaba,
Meguro-ku, Tokyo, 153-8505, Japan}
\author[3,4]{Toshiyuki Hirano}
\affil[4]{Tokyo Metropolitan College of Industrial Technology, 8-17-1, Minami-Senju, Arakawa-ku,
Tokyo, 116-8523, Japan}
\author[3]{Saisei Tahara}
\author[3]{Fumitoshi Sato}

\title{Time Evolution Simulation of a Hydrogen Atom Using Full Quantum Circuits}
% Use the \date command for email address(s) of corresponding authors
\date{*Email: hideo-t@iis.u-tokyo.ac.jp}

\begin{document}

\maketitle

\begin{abstract}

A quantum circuit for a first-quantization-based Hamiltonian simulator capable
of operating on a small-scale quantum computer emulator with fewer than 30
qubits was constructed.

To construct the circuit within limited quantum resources (number of qubits), we
applied Uniformly Controlled Rotation (UCR) circuits to calculate the Coulomb
potential. Furthermore, as a guideline for obtaining highly accurate
computational results, we established a method for handling the singularity
(pole) of the Coulomb potential and optimizing the discretization parameters.

Regarding the singularity of the Coulomb potential, we demonstrated that
divergence can be avoided and errors minimized by determining an effective value
from a simple geometric average of the potential within a minute circular region
centered on the pole.

Additionally, we proposed a method to determine the optimal spatial grid spacing
for a given number of qubits for coordinates, by evaluating the tradeoff between
the spatial discretization error caused by grid coarseness and the domain
truncation error due to a finite domain length. For the time evolution step
size, rather than relying on the formal Trotter error bound that assumes the
worst-case scenario, we clarified the conditions required to maintain sufficient
accuracy while avoiding an unnecessary increase in quantum circuit depth. This
was achieved based on the sampling theorem, focusing on the maximum frequency
component of the wave function.

\end{abstract}

%\section*{Keywords}
%Some journals require keywords: these normally should be given immediately
%after the abstract.

%\section*{Abbreviations}

%Some journals require a list of abbreviations: these normally should be given
%immediately after the keyswords (if required).

%%%%%%%%%%%%%%%%%%%%%%%%%%%%%%%%%%%%%%%%%%%%%%%%%%%%%%%%%%%%%%%%%%%%%
%% Start the main part of the manuscript here.
%%%%%%%%%%%%%%%%%%%%%%%%%%%%%%%%%%%%%%%%%%%%%%%%%%%%%%%%%%%%%%%%%%%%%
\subfile{jctc-draft_1_intro.tex}
\subfile{jctc-draft_2_method.tex}
\subfile{jctc-draft_34_results_discussions.tex}
% \subfile{jctc-draft_3_results.tex}
% \subfile{jctc-draft_4_analysis.tex}

%%%%%%%%%%%%%%%%%%%%%%%%%%%%%%%%%%%%%%%%%%%%%%%%%%%%%%%%%%%%%%%%%%%%%
%% The conclusion section
%%%%%%%%%%%%%%%%%%%%%%%%%%%%%%%%%%%%%%%%%%%%%%%%%%%%%%%%%%%%%%%%%%%%%

\section{Conclusion}
In this study, we presented a circuit configuration capable of simulating the
full time evolution of a two-dimensional hydrogen atom model on a Qiskit
emulator. To the best of our knowledge, this is the first report to implement
grid-based time evolution calculations, including the Coulomb potential,
entirely as a quantum circuit and to confirm its operation on an emulator.

Regarding the handling of the Coulomb potential singularity, we proposed a
method using $\reff=\deltar/4$ as the effective distance to determine the
effective potential at the grid point containing the singularity in the 2D
model. Numerical experiments confirmed not only that this value minimizes the
energy error of the system, but also that it coincides with the value derived
from the geometric area average around the singularity, clarifying the physical
validity of this approach in a discretized space.

Furthermore, we established a procedure for determining simulation parameters
that minimize errors under a given number of qubits ($\nb$). We demonstrated
that the spatial grid spacing, $\deltar$, can be determined as the optimal value
by balancing the tradeoff between spatial discretization error and domain
truncation error. This value determines the range starting at the ground state
up to the maximum quantum number that can be represented, which is the
representable dynamic range of the simulation.

Moreover, regarding the time step, $\deltat$, of the Trotter decomposition, we
numerically confirmed that it can be determined not from the formal upper bound
of the Trotter error assuming the worst case, but from the requirement of the
sampling theorem (half-period condition) based on the maximum frequency
component of the wave function. Particularly when targeting states near the
ground state, we demonstrated that within the range satisfying this condition,
the spatial error becomes dominant over the temporal error (i.e., the error
saturates), providing a guideline to avoid unnecessary increases in quantum
circuit depth.

These findings are expected to serve as highly practical and powerful guidelines
for designing and executing quantum simulations using the grid method under
limited quantum resources.

%\subsection{References}

\printbibliography[heading=bibintoc, title={References}]

\section*{Acknowledgements}

Please use ``The authors thank \ldots'' rather than ``The authors would like to
thank \ldots''.

\section*{Supporting information}

A listing of the contents of each file supplied as Supporting Information
should be included. For instructions on what should be included in the
Supporting Information as well as how to prepare this material for
publications, refer to the journal's Instructions for Authors.

The following files are available free of charge.
\begin{itemize}
  \item Filename-1: brief description
  \item Filename-2: brief description
\end{itemize}

\subfile{jctc-draft_a_appendix.tex}

\newpage

\rule{0.05in}{1.75in}%
\begin{minipage}[b][1.75in]{3.25in}
  \sffamily
  \frenchspacing

  Some journals require a graphical entry for the Table of Contents. This
  should be laid out ``print ready'' so that the sizing of the text is correct.

  The space available depends on the journal: J. Am. Chem. Soc. allows 3.25 in
  by 1.75 in and requires sanserif text. Some journals want different sizes:
  you can easily adjust here.
  
  The two rules either side of the content are there to help judge the height
  of your material: they may be deleted once not required.
  
\end{minipage}%
\rule{0.05in}{1.75in}

\end{document}
